\documentclass[11pt,a4paper]{article}
\usepackage{amsmath,amssymb,bm,geometry}
\usepackage[hidelinks]{hyperref}
\geometry{margin=2.2cm}
\newcommand{\C}{\mathbb{C}}
\newcommand{\E}{\mathbb{E}}
\newcommand{\rx}{\mathrm{rx}}
\newcommand{\tx}{\mathrm{tx}}
\newcommand{\tr}{\mathrm{tr}}
\title{Draft: Communication-Preserving Near-Field RIS Sensing via Phase Dithering}
\author{}
\date{}

\begin{document}
\maketitle

\section{Motivation and Objective}
This draft investigates a sensing mode where the RIS keeps its main focusing profile toward a served UE, while applying small phase dithers to extract environment information from CSI variations. The target is static-object detection (and optional localization) with minimal communication degradation.

\paragraph{Core question.}
Can repeated CSI under communication-preserving RIS micro-dithers provide enough information to detect a static object near the RIS in a near-field regime?

\paragraph{Main contributions.}
This draft makes six concrete contributions: (i) it introduces a communication-preserving sensing mode where RIS micro-dithers are applied around a UE-serving profile; (ii) it formulates a differential CSI model for static-object detection and optional localization; (iii) it derives a nuisance-aware EFIM/CRB feasibility framework with PEB interpretation; (iv) it presents a GLRT-based detection formulation with ROC-oriented evaluation; (v) it casts dither codebook design under explicit communication constraints such as SINR/BLER degradation limits; and (vi) it defines a practical evaluation protocol covering robustness to non-ideal RIS behavior and calibration drift.

\section{Related Work and Novelty Position}
Prior RIS codebook and ISAC studies mostly focus on one of the following: (i) communication beam training and link gain maximization, (ii) user localization support, or (iii) dual-purpose sensing/communication codebooks without explicitly preserving a fixed serving beam during sensing.

\paragraph{Representative literature.}
Recent works include learning-based RIS reflection codebooks for communication-oriented beam management, near-field ISAC codebooks for joint angle-range estimation, and near-field RIS localization studies based on FIM/PEB under practical phase-amplitude constraints.

\paragraph{Gap addressed in this draft.}
This work targets a different operating point: static-object sensing in near field while maintaining the RIS main profile for a served UE. Sensing is performed through small perturbations around the serving profile and differential CSI analysis.

\paragraph{Novelty dimensions.}
\begin{itemize}
\item \textbf{Problem setup:} communication-preserving static-object sensing, not generic beam training.
\item \textbf{Codebook design:} micro-dither codebooks constrained by communication-quality budgets.
\item \textbf{Theory:} unified EFIM/CRB (estimation feasibility) and GLRT/ROC (detection capability).
\item \textbf{Practicality:} explicit robustness analysis under RIS non-idealities and calibration drift.
\end{itemize}

\section{System and Signal Model}
Consider fixed BS, UE, and RIS positions over a sensing window. Let $m=0,\ldots,M-1$ index RIS states (with $m=0$ the baseline serving profile), $k=1,\ldots,K$ OFDM subcarriers, and $t=1,\ldots,T$ snapshots.

Let $\bm{\Gamma}_0 \in \C^N$ be the baseline RIS profile focused on UE. Dithered profiles are
\begin{equation}
\bm{\Gamma}_m = \bm{\Gamma}_0 \odot e^{j\bm{\delta}_m},
\end{equation}
where $\bm{\delta}_m$ is a small phase perturbation vector.

The received CSI sample is modeled as
\begin{equation}
y_{m,k,t} = \mu_{m,k}(\bm{\theta}) + n_{m,k,t}, \quad n_{m,k,t}\sim\mathcal{CN}(0,\sigma^2),
\end{equation}
with
\begin{equation}
\mu_{m,k}(\bm{\theta}) =
h^{\mathrm{dir}}_k + h^{\mathrm{RIS-UE}}_{m,k}
+ h^{\mathrm{obj}}_{m,k}(\bm{\theta}_o) + h^{\mathrm{nu}}_{m,k}(\bm{\eta}).
\end{equation}
Here, $\bm{\theta}_o=[x_o,y_o,z_o,\alpha_o]^T$ are object parameters (position and complex reflectivity), and $\bm{\eta}$ denotes nuisance parameters (residual phase offsets, calibration drift, etc.).

\section{Differential Observation for Robustness}
To suppress common terms and hardware drifts, use differential CSI relative to baseline:
\begin{equation}
\Delta y_{m,k,t}=y_{m,k,t}-y_{0,k,t},\quad m\ge 1.
\end{equation}
For static BS/UE/RIS geometry, $\Delta y_{m,k,t}$ enhances object-dependent modulation induced by RIS dithers.

\section{Feasibility via FIM and EFIM}
Let $\bm{\theta}=[\bm{\theta}_o^T,\bm{\eta}^T]^T$. Under complex Gaussian noise, the Fisher Information Matrix is
\begin{equation}
\mathbf{J}(\bm{\theta})=
\sum_{m,k,t}\frac{2}{\sigma^2}
\Re\!\left\{
\left(\frac{\partial \mu_{m,k}}{\partial \bm{\theta}}\right)^H
\left(\frac{\partial \mu_{m,k}}{\partial \bm{\theta}}\right)
\right\}.
\end{equation}

Partition
\begin{equation}
\mathbf{J}=
\begin{bmatrix}
\mathbf{J}_{oo} & \mathbf{J}_{o\eta}\\
\mathbf{J}_{\eta o} & \mathbf{J}_{\eta\eta}
\end{bmatrix}.
\end{equation}
The equivalent Fisher matrix for object parameters is
\begin{equation}
\mathbf{J}_{\mathrm{eq}}=\mathbf{J}_{oo}-\mathbf{J}_{o\eta}\mathbf{J}_{\eta\eta}^{-1}\mathbf{J}_{\eta o}.
\end{equation}
Then CRB is
\begin{equation}
\mathrm{cov}(\hat{\bm{\theta}}_o)\succeq \mathbf{J}_{\mathrm{eq}}^{-1}.
\end{equation}
If estimating 3D position, the position error bound (PEB) is
\begin{equation}
\mathrm{PEB}=\sqrt{\tr\!\left(\left[\mathbf{J}_{\mathrm{eq}}^{-1}\right]_{1:3,1:3}\right)}.
\end{equation}

\paragraph{Interpretation.}
This bound quantifies when micro-dithering is informative enough for localization and how much near-field geometry and dither design improve observability.

\subsection{When object localization is identifiable}
Localization is possible only if the measurement model is sufficiently informative and non-degenerate. In practice, the following conditions should hold:
\begin{itemize}
\item \textbf{Sufficient diversity across dithers:} the Jacobian vectors $\partial \mu_{m,k}/\partial \bm{\theta}_o$ should not be collinear across $m$.
\item \textbf{Frequency diversity:} enough subcarriers $K$ and bandwidth are needed to separate range-sensitive effects.
\item \textbf{Geometry diversity:} RIS-BS-UE-object geometry should avoid near-symmetric layouts that create ambiguous solutions.
\item \textbf{Near-field modeling:} spherical-wave path terms must be used when the object is in Fresnel region.
\item \textbf{Nuisance control:} phase/CFO/calibration drifts must be mitigated (or marginalized through EFIM).
\end{itemize}
If these conditions are weak, detection can still be feasible while localization error remains large.

\subsection{How to report CRB results}
For each candidate setup (SNR, $M$, $K$, RIS size, object distance), report:
\begin{itemize}
\item $\lambda_{\min}(\mathbf{J}_{\mathrm{eq}})$ as an observability indicator;
\item PEB from $\mathbf{J}_{\mathrm{eq}}^{-1}$;
\item empirical localization RMSE of the estimator;
\item the ratio RMSE/PEB to quantify estimator efficiency.
\end{itemize}
This directly answers whether localization is feasible in each operating region.

\section{Detection Formulation (Primary Task)}
For static-object presence testing:
\begin{align}
\mathcal{H}_0 &: \Delta y_{m,k,t} = w_{m,k,t},\\
\mathcal{H}_1 &: \Delta y_{m,k,t} = s_{m,k}(\bm{\theta}_o)+w_{m,k,t}.
\end{align}
A generalized likelihood ratio test (GLRT) is
\begin{equation}
\Lambda(\Delta\mathbf{y})=
\frac{\max_{\bm{\theta}_o,\bm{\eta},\sigma^2} p(\Delta\mathbf{y}\mid\mathcal{H}_1)}
{\max_{\bm{\eta},\sigma^2} p(\Delta\mathbf{y}\mid\mathcal{H}_0)}
\gtrless_{\mathcal{H}_0}^{\mathcal{H}_1}\gamma.
\end{equation}
Report ROC curves, AUC, and $P_D$ at fixed $P_{FA}$.

\section{Dither Codebook Design Under Communication Constraint}
Design dithers to maximize sensing information while preserving UE link:
\begin{align}
\max_{\{\bm{\delta}_m\}} \quad & \log\det\!\big(\mathbf{J}_{\mathrm{eq}}(\{\bm{\delta}_m\})\big) \\
\text{s.t.}\quad
& \Delta\mathrm{SINR}(\{\bm{\delta}_m\}) \le \epsilon, \\
& |\delta_{m,n}| \le \delta_{\max},\quad \forall m,n.
\end{align}
Alternative criteria include minimizing $\tr(\mathbf{J}_{\mathrm{eq}}^{-1})$ or maximizing deflection coefficient for detection.

\section{Evaluation Plan}
\begin{itemize}
\item \textbf{Setup:} fixed BS/UE/RIS, near-field RIS model, static objects at known positions.
\item \textbf{Data:} collect CSI over $M$ dither states, $K$ subcarriers, $T$ repeats per state.
\item \textbf{Tasks:} object detection; optional coarse object localization.
\item \textbf{Metrics:} ROC/AUC/$P_D$, localization RMSE vs CRB/PEB, SINR/BLER/throughput impact, sensing overhead.
\item \textbf{Robustness:} cross-session and calibration-drift tests.
\end{itemize}

\subsection{Simulation checklist for localization feasibility}
\begin{enumerate}
\item \textbf{Baseline run:} no-object and object-present datasets with fixed geometry.
\item \textbf{SNR sweep:} at least 4--6 SNR points to identify threshold behavior.
\item \textbf{Dither sweep:} vary number of states $M$ and dither amplitude $\delta_{\max}$.
\item \textbf{Bandwidth sweep:} vary $K$ and effective bandwidth to test range sensitivity.
\item \textbf{Geometry sweep:} object position grid (distance and azimuth/elevation sectors).
\item \textbf{Nuisance sweep:} inject phase drift/CFO/calibration errors and compare FIM vs EFIM predictions.
\item \textbf{Communication impact:} report SINR, BLER, throughput drop during sensing.
\item \textbf{Bound comparison:} plot RMSE against PEB for each sweep.
\end{enumerate}

\subsection{Simulation methodology (implementation-oriented)}
The following procedure can be implemented directly in a Monte Carlo simulator:
\begin{enumerate}
\item \textbf{Initialize geometry and channels:} set fixed BS/UE/RIS positions, near-field RIS model, and an object-location grid. Generate channels for direct, RIS-assisted serving, and object-scattered components.
\item \textbf{Construct RIS states:} compute a serving profile $\bm{\Gamma}_0$ for UE communication, then generate $M-1$ micro-dither profiles $\{\bm{\Gamma}_m\}_{m=1}^{M-1}$.
\item \textbf{Generate CSI measurements:} for each hypothesis ($\mathcal{H}_0$ no object, $\mathcal{H}_1$ object present), each RIS state $m$, subcarrier $k$, and repeat $t$, synthesize $y_{m,k,t}$ under AWGN and selected non-idealities.
\item \textbf{Apply differential preprocessing:} compute $\Delta y_{m,k,t}=y_{m,k,t}-y_{0,k,t}$ and extract detection/localization features.
\item \textbf{Run detector and estimator:} evaluate GLRT (or baseline detector) for ROC/AUC and run localization estimator when $\mathcal{H}_1$ is true.
\item \textbf{Compute bounds and communication impact:} evaluate EFIM/PEB for each setup and report SINR/BLER/throughput changes during sensing.
\item \textbf{Aggregate over trials:} run sufficient Monte Carlo trials and summarize mean and confidence intervals.
\end{enumerate}

\subsection{Default parameter ranges}
Unless otherwise specified, a practical initial sweep is:
\begin{itemize}
\item \textbf{SNR:} $[-10, -5, 0, 5, 10, 15, 20]$ dB.
\item \textbf{Dither states:} $M \in \{4, 8, 16, 32\}$.
\item \textbf{Dither amplitude:} $\delta_{\max} \in \{5^\circ, 10^\circ, 20^\circ, 30^\circ\}$.
\item \textbf{Subcarriers:} $K \in \{32, 64, 128, 256\}$.
\item \textbf{Repeats per state:} $T \in \{4, 8, 16\}$.
\item \textbf{Monte Carlo trials:} at least 200 per setup (increase to 500 for final plots).
\end{itemize}
These ranges can be tightened once a feasible operating region is identified.

\section{Expected Contributions}
\begin{enumerate}
\item Communication-preserving near-field RIS sensing via micro-dithers around serving beam.
\item Unified estimation and detection feasibility analysis using EFIM/CRB and GLRT.
\item Dither codebook design with explicit communication-quality constraints.
\item Practical validation showing detection/localization gains at low communication overhead.
\end{enumerate}

\section{Conclusion}
The proposed framework extends opportunistic RIS-ISAC to static-object sensing without abandoning UE-oriented beam focusing. The feasibility analysis (EFIM/CRB) and detection theory (GLRT/ROC) provide a rigorous path to justify and optimize the approach.

\end{document}
